\section{Giới thiệu chung}

\subsection{Giới thiệu công ty}

KiotViet là thương hiệu phần mềm quản lý bán hàng do Công ty Cổ phần Công nghệ Citigo phát triển, ra đời với mục tiêu giúp những cửa hàng vừa và nhỏ ở Việt Nam quản lý việc bán hàng bằng phần mềm thay vì bằng sổ tay. Từ nhóm khách hàng ban đầu là các cửa hàng bán lẻ, công ty mở rộng dần sang nhà hàng, quán cà phê, cửa hàng làm đẹp và khách sạn; mỗi ngành có cách vận hành riêng nên được tách thành những dòng sản phẩm khác nhau dùng chung phần nền bên dưới.

Em thực tập tại nhóm phát triển sản phẩm cho ngành khách sạn. Sản phẩm phục vụ khách sạn, nhà nghỉ, homestay và resort, với người dùng cuối là lễ tân, thu ngân, buồng phòng, kế toán và chủ cơ sở lưu trú. Điểm đáng chú ý là sản phẩm không chỉ chạy ở Việt Nam mà còn được dùng ở một số nước khác trong khu vực, nên phần mềm phải nói được nhiều thứ tiếng và tính toán đúng theo múi giờ, đơn vị tiền tệ của từng nơi.

\subsection{Giới thiệu công việc}

Em tham gia nhóm với vai trò thực tập sinh phát triển phần mềm, làm cả phần giao diện và phần xử lý phía máy chủ, từ ngày 12 tháng 3 năm 2026 đến ngày 7 tháng 9 năm 2026.

Ngày đầu em hơi choáng. Sản phẩm không nằm trong một kho mã duy nhất mà chia thành chín kho mã riêng: màn hình cho lễ tân, màn hình cho người quản lý, dịch vụ xử lý đơn đặt phòng, dịch vụ sinh báo cáo, dịch vụ xuất nhập tệp dữ liệu, phần giao tiếp với ứng dụng di động và phần nền dùng chung. Tổng cộng có hàng trăm nghìn dòng mã đã tồn tại từ trước, do rất nhiều người viết qua nhiều năm. Em mất gần hai tuần đầu chỉ để trả lời được câu hỏi đơn giản nhất: khi lễ tân bấm một nút trên màn hình thì dữ liệu chạy qua những đâu.

Công việc của em chia thành bốn nhóm: sửa lỗi được báo về từ khách hàng đang dùng thật, làm tính năng mới theo yêu cầu, tối ưu hiệu năng ở những chỗ làm hệ thống chậm hoặc chặn nhau, và tham gia dựng phần quản lý phiên bản mới.

\begin{table}[htbp]
\centering
\caption{Tóm tắt đóng góp trong kỳ thực tập}
\label{bang:dongghop}
\small
\setstretch{1.1}
\begin{tabular}{@{}>{\raggedright\arraybackslash}p{7.6cm}r@{}}
\toprule
\textbf{Hạng mục} & \textbf{Số liệu} \\
\midrule
Thời gian thực tập & 12/03/2026 đến 07/09/2026 \\
Số kho mã đã tham gia & 9 \\
Số lần ghi mã, không tính lần gộp nhánh & 698 \\
Số đầu việc khác nhau đã xử lý & 342 \\
Số phiếu công việc đã hoàn thành & 173 \\
Số ngày có ghi mã & 99 \\
Số dòng mã thêm mới & 34.994 \\
Số dòng mã xoá bỏ & 10.079 \\
Số lượt sửa tệp & 1.947 \\
Số dòng mã kiểm thử tự động đã viết & 1.919 \\
\bottomrule
\end{tabular}
\end{table}

Bảng \ref{bang:dongghop} tóm tắt phần đóng góp của em, số liệu lấy trực tiếp từ lịch sử mã nguồn. Con số dòng mã đã loại bỏ các tệp bản dịch ngôn ngữ, tệp cấu hình sinh tự động và tệp nhị phân, để phản ánh đúng phần mã do em thực sự viết.

\subsection{Giới thiệu qua bài toán}

Nghe tên thì phần mềm quản lý khách sạn có vẻ đơn giản: khách đến thì nhận phòng, khách đi thì trả phòng và thanh toán. Nhưng khi bắt tay vào làm em mới thấy phần khó không nằm ở chỗ lưu dữ liệu, mà ở chỗ cùng một dữ liệu phải trả lời đúng cho rất nhiều câu hỏi khác nhau, ở nhiều nơi khác nhau, tại nhiều thời điểm khác nhau.

Lấy một ví dụ em gặp thật. Một phòng bị khoá lại để sửa chữa trong ba ngày. Việc này nghe rất nhỏ, nhưng ngay sau đó hàng loạt câu hỏi khác phải đổi câu trả lời: màn hình lễ tân phải không cho đặt phòng đó nữa, lịch phòng phải vẽ ra một khoảng bị chặn, báo cáo công suất phòng phải hiểu rằng phòng này đang không kinh doanh nên không được tính vào mẫu số, và nếu người quản lý xem lại số liệu tháng trước thì phòng đó lại phải được tính bình thường vì lúc ấy nó chưa bị khoá. Một thay đổi ở một chỗ kéo theo hệ quả ở năm sáu chỗ khác, và những chỗ đó do người khác viết từ trước.

Bài toán còn khó thêm vì ba lý do. Thứ nhất, hệ thống đang chạy thật nên không có chuyện làm lại từ đầu. Thứ hai, một khách hàng có thể có nhiều chi nhánh ở nhiều nơi, mỗi nơi một múi giờ và một đồng tiền, nên câu hỏi tưởng dễ như hôm nay bán được bao nhiêu cũng cần định nghĩa rõ hôm nay là theo giờ ở đâu. Thứ ba, phần mềm hiển thị được năm ngôn ngữ, nên mỗi chuỗi chữ, mỗi định dạng ngày và mỗi cách viết số đều phải xử lý cho cả năm. Trong bối cảnh đó, phần việc em nhận không phải là xây một hệ thống mới, mà là bổ sung và làm tốt hơn những phần cụ thể của một hệ thống đang sống.
